\documentclass[11pt]{article}
\usepackage{preprint}
\newfontfamily\ipafont{LibertinusSerif-Regular.otf}

\title{Enhancing Writing Skills for Chinese Students in Grades 7--9}
% authblk prints only the last \author when no \affil is given:
% list every author on one line instead (no affiliation marks).
\makeatletter
\renewcommand\AB@authnote[1]{}
\renewcommand\@author{\AB@authlist}
\makeatother
\author{Yuning Li}
\author{Jiadi He\email{jh2288@georgetown.edu}}
\author{Ningxi Deng}
\shorttitle{Enhancing Writing Skills, Grades 7--9}
\runningauthors{Li, He \& Deng}
\keywords{learning design; language education; educational technology}

\begin{document}
\maketitle

\begin{abstract}
This project conducted an in-depth investigation of Chinese middle school students' competence in English writing, particularly in grades seven through nine, and assessed the effectiveness of current English education methods in meeting students' writing needs. The study points out that although students already possess basic language skills, the lack of sufficient writing practice, limited practical application scenarios, and the monotonous nature of traditional teaching methods have not been effective in stimulating students' interest and creativity in writing, making it difficult for them to express complex ideas and emotions in writing.
\end{abstract}

\section{Problem/Crisis Identified}

The purpose of this project is to identify and analyze the shortcomings of current English writing teaching strategies, especially the lack of engaging classroom content and the failure of teaching methods to adequately adapt to students' specific needs. In this regard, the study proposes to improve students' English expressive skills by revamping the teaching methodology to incorporate more varied and practical writing tasks, as well as frequent and targeted feedback and guidance.

By improving curriculum design and pedagogical practices, this project aims to provide educators with concrete strategies and recommendations to make English writing instruction more effective and responsive to students' future academic and professional needs. We expect that through these improvements, students will not only enhance their writing skills, but also build greater self-confidence in the learning process and be better prepared to meet future challenges, thus finding their place in a globalized language environment.

In a recent analysis of students' writing samples from the class's official WeChat account, a phenomenon emerged concerning the topic of ``Travel Experience in Shandong.'' The students' essays focused on their travels to Mount Tai, written in the style of letters to friends or family.

The students' responses to the assignment were notably articulate, showing a deep appreciation for the cultural and natural richness of Mountain Tai. For example, one student's reflection on the enduring natural beauty of Shandong, spanning ancient to modern times and symbolizing a collective heritage, added profound depth to their narrative. This perspective portrayed Shandong as not just a place but a ``natural symbol of ancient China.''

However, alongside this richness in descriptions, the students' writings often exhibited clichés and repetitions regarding their experiences, particularly in their descriptive elements. Many students tended to use common phrases or clichéd expressions when describing their travels to Mount Tai in Shandong. For example, phrases like ``breathtaking views'' and ``serene landscapes'' were frequently employed without delving into specific details or unique perspectives.

Moreover, there was a noticeable pattern of similar experiences being narrated across different essays. Multiple students described similar scenes of hiking up Mount Tai, exploring ancient temples, and expressing feelings of awe and admiration at the natural beauty of the surroundings. While these experiences are undoubtedly significant and impactful, the lack of distinctiveness in how they were portrayed indicates a gap in students' writing in terms of originality and creativity. Instead of offering unique perspectives or personal anecdotes that could have added depth to their narratives, many students relied on generic descriptions and common sentiments. This trend highlighted a missed opportunity for students to showcase their individuality and creativity in storytelling. For instance, rather than simply stating admiration for the mountain's grandeur, students could have delved into specific moments that stood out to them, such as interactions with locals, unexpected challenges during the hike, or insights gained from visiting historical sites.

\section{Analyze the Problem/Crisis}

The analysis of student essays on the theme ``Shandong Travel Experience'' revealed a series of challenges in writing training that reflect deeper educational issues. Lack of an English immersion environment, limitations of traditional teaching methods, limited feedback and guidance, cultural and linguistic barriers, and inadequate technology integration were the main underlying factors.

Firstly, students often have difficulty practicing English in authentic contexts because of the lack of natural English immersion environments, which is especially evident when they are describing complex emotions and specific scenarios in English. This lack of environment leads students to rely too much on standardized phrases and expressions from textbooks and lack the ability to create new expressions and descriptions, thus making their compositions seem repetitive and uninspired.

Secondly, traditional, textbook-dependent teaching methods are often ineffective in engaging students and stimulating their interest and creativity in writing. This approach usually ignores students' individualized learning needs and the development of creative thinking, so that students tend to focus on the correct use of grammar and vocabulary in their writing at the expense of innovation in content and the presentation of their individual voices.

In addition, students do not receive targeted feedback and guidance during the writing process, which limits their opportunities to improve their writing skills and develop a personal style of expression. Effective feedback should include specific suggestions for improvement and elements of encouragement to help students understand their strengths and areas for improvement. Schools may consider using more peer review and one-on-one tutoring with teachers to provide more detailed and personalized writing instruction.

At the same time, students often encounter cultural and linguistic barriers to expressing themselves in English due to the large differences between English and their native language and cultural backgrounds. This not only limits their ability to express themselves in writing, but also affects their comprehension and assimilation of English materials. To help students overcome these barriers, educators can introduce more instructional content about cultural contexts, such as through cultural exchange programs or bringing in foreign teachers to increase students' direct exposure to and understanding of the target language and culture.

Finally, although digital tools have begun to be widely used in teaching and learning, they often lack full technological integration to provide interactive and real-time learning experiences. This limits the opportunities for students to practice and innovate in language through technology.

Of these essential factors, traditional teaching methods are particularly critical because they directly impact how students engage with the material they are learning and profoundly affect their ability to develop critical thinking and creative writing skills.

Traditional teaching methods often rely on textbooks and teacher-centered instruction, an approach that often fails to engage students on a deeper, more interactive level. This is critical for meaningful learning and mastery of complex skills such as writing. When students are not fully engaged in the learning process, their motivation and interest in learning diminishes, which in turn affects writing achievement and creativity. A recent study emphasized that ESL writing anxiety is largely influenced by traditional and less interactive teaching methods, which can exacerbate students' anxiety and reduce their ability to complete writing tasks. (Wern \& Rahmat, 2021).

In addition, traditional teaching methods tend to discourage personal expression or creativity, both of which are crucial for effective writing in any language, especially a second language like English. Personal expression and creativity are essential elements of writing that help students develop a unique voice and perspective in their writing, thereby enhancing the quality and individuality of their writing.

Addressing the limitations of traditional teaching methods not only increases student engagement and creativity, but also provides important opportunities for innovation in teaching practices. By introducing more interactive, student-centered learning methods, educators can better meet the diverse needs of learners and promote active learning and critical thinking. This may include project-based learning, collaborative learning, and the use of technology tools for interactive exploration and practice. Research also suggests that incorporating affective strategies and active error correction can help reduce writing anxiety, which is critical to improving the learning environment. (Wern \& Rahmat, 2021)

In summary, focusing on and improving traditional teaching methods to include more interactive and student-centered teaching strategies will not only improve students' English writing skills, but also stimulate their creativity and critical thinking, laying the foundation for better expression and communication in a globalized context.

\section{Product}

This English writing program for Chinese middle school students (especially those in grades 7--9) is designed to enhance students' writing skills in a holistic manner through structured learning tasks and creative, interactive activities. The course is broken down into four consecutive one-week phases, each of which focuses on strengthening students' skills in different aspects of the writing process, from interpreting writing prompts to final revisions of their essays.

\subsection{Sample Activity}

During the first week of the course, we focus on teaching students how to understand and analyze writing prompts, a skill they must master before they begin writing. Through interactive digital slides, students are guided through an in-depth exploration of a variety of writing topics and corresponding great examples. These slides not only demonstrate how to respond effectively to writing prompts, but also include detailed analysis of the model essays so that students can see examples of high-quality writing. In small groups, students take on different roles, such as author or reviewer, and such role-playing activities allow students to learn by doing how to critique and discuss essays, improving their analytical and critical thinking skills. For example, an activity might consist of the instructor presenting a writing prompt for an expository essay and then asking students to analyze a sample essay in response to this prompt, discussing the validity of its thesis and the adequacy of its arguments.

Moving into the second week, the course will turn to how to organize ideas and construct an outline, a critical step in the writing process. Students will use mind-mapping software and collaborative brainstorming tools that help students organize and link their ideas in a visual way. In group discussion sessions, students will work together to create mind maps around a common essay topic, such as climate change. Each student will be required to contribute at least two unique ideas or arguments and explore how these ideas can be logically integrated into the essay. This process not only fosters collaboration among students, but also helps them expand their thinking and develop a more sophisticated essay structure through collective wisdom.

In the third week, students will begin drafting specific essays based on the outlines they have developed. By utilizing real-time language feedback tools such as Grammarly, students receive instant feedback on grammar, spelling, and style as they write, which can help them identify and correct errors in a timely manner. Additionally, the week includes a specialized writing workshop in which the instructor demonstrates how to avoid common logical fallacies and shows how to use rhetorical techniques to enhance the persuasive power of an essay, such as comparisons and analogies.

The final week of the course focuses on essay revision and peer review. By using a peer review platform such as Peergrade, students not only learn how to give constructive feedback to others, but also receive comments from their peers on their own essays. This stage of interaction helps students understand the quality of writing from different perspectives and promotes deeper reflection on their own work. Teachers have a crucial role at this stage, not only overseeing the review process but also providing expert feedback to help students identify and improve weaknesses in their essays.

By designing the program in such a coherent and integrated way, we aim to not only enhance students' writing skills, but also, through continuous interaction and feedback, to increase their self-confidence and critical thinking skills, thus laying a solid foundation for their future academic and professional careers.

In this innovative English writing program, a series of activities have been carefully designed to ensure that students are able to improve their writing skills in a variety of ways. Through these activities, students will engage in a rich learning experience that includes practice in all aspects of listening, speaking, reading and writing, aiming to improve their overall language use.

\subsection{Learner Engaged Activity}

Students will engage in a variety of interactive activities that include role-plays, group discussions, real-time writing exercises, and peer reviews. In role-plays, students will simulate different writing scenarios, for example, playing the role of a news reporter or a critic, analyzing and discussing different perspectives on a given topic. Such activities help students practice how to use appropriate language and style in a given situation. In group discussions, students will work together to explore how to construct an argument, share their ideas and receive feedback from their peers. These discussions not only enhance students' critical thinking but also improve their communication skills.

Real-time writing exercises allow students to improve their drafts with real-time guidance from the instructor. By using writing software with grammar and style checking features, students are able to receive instant feedback and adjust their texts. Finally, through peer review, students will learn how to objectively evaluate the work of others and provide constructive feedback, as well as learn how to improve their own writing from the feedback of others. According to Graham and Perin (2007), developing instructional arrangements in which adolescents work together to plan, draft, revise, and edit their compositions could enhance students’ writing skills. With high levels of interactions with peers, Students will reflect on the feedback received, make thoughtful revisions to their drafts, and apply the insights gained from the peer review process.

\subsection{Reasons for selecting these components/activities}

These activities have been chosen and structured in this way to ensure that students are exposed to and practice writing in English from multiple perspectives while improving their analytical and critical thinking skills. Through role-playing, students are able to apply their language skills in real-world simulated situations, an approach that enlivens the learning process. Group discussions and peer reviews, on the other hand, emphasize the importance of collaborative learning, and through sharing and feedback, students can better understand how to effectively communicate their ideas and learn from the feedback of others. Real-time writing exercises and immediate feedback mechanisms ensure that students can immediately correct their mistakes and deepen their understanding of language rules.

\subsection{How to connect learners to practical functions}

This course connects learners to the practical functions of writing skills by integrating real-world application of writing tasks and modern technological tools. For example, the use of real-time feedback tools not only helps students correct grammatical errors, but also demonstrates how to optimize sentence structure and enhance the logical flow and persuasiveness of a text through concrete examples. By simulating real-world writing tasks, such as writing a report or essay, students are able to apply the skills they learn in the course to prepare for the writing demands of their future academic and professional careers.

\subsection{Examples of interactions created}

In a typical interaction, the instructor might present a writing challenge such as, ``Suppose you are a reporter for a technology magazine and need to write an article about the latest cell phone technology. Please consider that your audience is tech enthusiasts.'' Students, in turn, will need to conceptualize and draft an article, during which the instructor will provide targeted questions to guide students' thinking, such as, ``Does your article contain key information that would influence a tech enthusiast's decision?'' Students will also interact with each other, and they may ask each other, ``Which description of cell phone technology do you think would best capture the reader's attention? Why?''

\subsection{Responding to ``I don't know'' situations}

In this English writing course, we place great emphasis on responding to the ``I don't know'' situations that students may encounter during the course of their studies. The course is designed with a particular focus on helping students overcome these challenges through concrete examples and practical exercises to ensure that they are able to learn and apply new vocabulary and language structures effectively.

\subsubsection*{Tackling the challenge of learning new words}

When students encounter unfamiliar words, such as `sustain', in their writing or reading, our digital learning platform automatically provides a range of resources to help with comprehension and memorization. These resources include the word's definition, pronunciation, and lexical properties, as well as a selection of example sentences designed to show how the word is used in different contexts. For example, for the word ``sustain'', the system might display the following information:

\begin{quote}
\begin{minipage}{\linewidth}
Definition: to hold for a long time; support, sustain. \\
Pronunciation: {\ipafont /səˈsteɪn/}. \\
Example sentence:
\begin{itemize}[label={},leftmargin=1.5em,topsep=0pt,itemsep=0pt]
  \item Economic growth needs to sustain good policy support.
  \item They discuss how to SUSTAIN natural resources for use by future generations.
\end{itemize}
\end{minipage}
\end{quote}

In this way, students not only understand the meaning of the words, but also see how they are used in real-life utterances, which helps them use these new vocabulary words more accurately in their own writing.

For more complex language structures, such as inverted sentences, the curriculum is designed to include teacher-led mini-lessons. In these lessons, teachers will first explain the definition and usage of inverted sentences, and then show how to use this structure in writing through specific examples. For example, a teacher might explain the difference between normal and inverted order and show the following examples:

\begin{quote}
Normal Order: Only then did he realize he was wrong. \\
Inversion: Only then he realized he was wrong.
\end{quote}

The teacher will then direct students to try to apply these structures to their own writing exercises, possibly setting up a mini-exercise such as, ``Try to write a sentence describing a sunrise using inverted clauses.'' Students' attempts might be, ``Only at dawn did the sun peek over the horizon.''

Through this teaching method, students not only understand the construction of inverted sentences, but also master their use by practicing them themselves, and this combination of theory and practice greatly enhances learning.

\subsection{Feedback mechanisms}

In this English writing program, a multi-level feedback mechanism is used to ensure that students learn from their mistakes and improve their writing skills in response to the errors they may encounter in the writing process. This mechanism involves immediate technical feedback, peer review, and professional guidance from the instructor, each of which is designed to help students understand and correct their errors while providing specific suggestions to avoid future mistakes.

When students use incorrect tenses or inflections in their writing, such as incorrectly using the past tense as the present tense, the writing assistance software we use, such as Grammarly, will immediately recognize these errors. Not only will the software flag the error, but it will also provide the correct form of the inflection or tense and give an explanation as to why the original use was incorrect. For example, if a student writes, ``He goes to school every day,'' the software will point out the error, suggesting that it be changed to ``He goes to school every day,'' and explaining that this is because the subject is in the third person singular and the verb should be in the third person singular present tense.

During the peer review phase of an article, students review each other's manuscripts and provide feedback. This process encourages students to think critically and help each other improve their essays. For example, if a student's essay makes jumps in logic, a peer might comment, ``Your argument jumps around a bit in the second paragraph; try adding some transitional statements to connect the preceding and following text.'' This feedback prompts the writer to reflect on how to make the essay flow and learn how to enhance its coherence through logical connections.

Following the peer review, the teacher will provide more in-depth feedback by synthesizing the students' work. The teacher will point out logical weaknesses in the essay and provide specific suggestions for improvement. For example, in response to an illogical essay, the teacher might suggest, ``When discussing this idea, start by presenting some background information or relevant data, which will not only help the reader understand your position, but also enhance the persuasiveness of your argument.'' With such guidance, students are able to learn how to construct more rigorous and persuasive arguments.

Such a comprehensive feedback mechanism ensures that students receive ongoing support and guidance in the writing process, not only helping them to correct specific errors, but also teaching them how to think about and implement more effective writing strategies so that they can avoid the same mistakes in their future writing and improve their overall writing.

This comprehensive instructional design ensures the holistic development of students in the writing learning process through well-planned writing activities and a multi-level feedback system. This approach not only focuses on improving students' basic language skills such as grammar, vocabulary, and sentence structure, but also broadly emphasizes the importance of critical thinking, creativity, and effective communication, all of which are indispensable skills that students will need to face future academic and professional challenges.

The development of the Interactive English Writing Curriculum for Chinese middle school students presents several key advantages that align with the modern educational needs of a diverse student population. Firstly, the structured yet flexible design of the curriculum caters to different learning styles and paces, making it broadly accessible. The integration of digital tools and interactive activities like role-playing and peer reviews enhances student engagement, which is critical for sustained learning. This approach also allows for real-time feedback, a crucial element in language learning, helping students immediately understand and correct their errors, thereby accelerating the learning process.

Another significant advantage is the curriculum's emphasis on collaborative learning. By encouraging peer interactions through group discussions and reviews, students not only learn from the instructional content but also gain valuable insights from their peers. This social aspect of learning is vital for developing communication skills and can increase motivation and enjoyment, key factors in language acquisition.

The use of technology in this curriculum also prepares students for a globalized digital world. Familiarity with digital writing tools and online communication platforms equips them with skills that are increasingly demanded in higher education and the workplace. Moreover, the curriculum's focus on pragmatic language functions ensures that students are not just learning English in abstract but are applying it in real-world contexts, enhancing their ability to think critically and creatively.

Despite these strengths, the curriculum faces certain limitations and boundaries. One of the primary challenges is the reliance on digital tools, which assumes that all students have equal access to technology and the internet. In regions where this is not the case, the effectiveness of the curriculum could be compromised. Additionally, while the curriculum aims to be inclusive, the varying degrees of prior English exposure and proficiency among students can lead to disparities in how students engage with and benefit from the curriculum. Advanced students might find the pace too slow, while beginners could struggle with some of the more challenging activities.

Another limitation is the curriculum’s intensive use of peer feedback, which depends heavily on the students' ability to provide constructive criticism. This requires careful guidance and training, which can be time-consuming and challenging to implement effectively across diverse classrooms. Also, the teacher’s role is critical in facilitating discussions and ensuring that feedback is appropriate and helpful, placing a significant demand on the teacher’s time and skills.

\section{Processes of Eliciting and Selecting Feedback}

Feedback is a cornerstone of this curriculum, and its efficacy hinges on the robustness of feedback mechanisms. During the pilot phase, feedback was systematically elicited through surveys and focus groups involving students, teachers, and educational experts. Classroom observations and digital analytics provided additional insights into how well the activities and tools were being used and their impact on learning.

Selecting which feedback to incorporate involved prioritizing changes that could broadly enhance accessibility and engagement. Feedback that addressed technical issues with digital tools, clarity of instructions, and the effectiveness of peer interactions was prioritized to refine the curriculum further.

\section{Potential for Future Steps}

Looking ahead, there is substantial potential for expanding the curriculum to include more adaptive learning technologies that can tailor content to individual student needs, thereby addressing the challenge of varying proficiency levels. Another future step could involve broader collaboration with educational technology firms to enhance the digital platform’s capabilities, such as using AI to provide more nuanced feedback on student writing.

Additionally, expanding the curriculum to include other aspects of English, such as speaking and listening, could provide a more comprehensive language learning experience. Collaborations with schools in other regions could also be explored to adapt the curriculum for different cultural and educational contexts, enhancing its universality and relevance globally.

As the Interactive English Writing Curriculum for Chinese middle school students culminates, it's pivotal to reflect on the transformative educational journey it promises. Designed meticulously to address the core challenges faced by students in mastering English writing, this curriculum embodies a progressive approach to language education that harmonizes technology, creativity, and interaction.

\section{Central Takeaways}

\begin{enumerate}
  \item Comprehensive Integration of Technology and Interactive Learning: A standout feature of the curriculum is its adept use of digital tools to create a dynamic learning environment. This integration is crucial, not just for engaging students, but for preparing them for the technological demands of the contemporary world. The curriculum employs technology to provide instant feedback, helping students swiftly correct their errors and apply new knowledge in real-time, which is a significant advancement over traditional methods.
  \item Emphasis on Collaborative Learning: The curriculum fosters an educational atmosphere that is enriched by peer interactions through discussions and reviews, crucial for developing communication skills and enhancing student motivation. This collaborative framework is vital for nurturing a supportive learning community where students benefit from diverse perspectives and collective wisdom.
  \item Pragmatic Language Use: ``By aligning activities with real-world applications, the curriculum ensures that students are not merely learning English but are applying it practically.'' (Yang, 2013) This approach significantly boosts their ability to think critically and creatively, enhancing their capacity to communicate complex ideas effectively.
  \item Adaptability and Future Prospects: ``The flexible design of the curriculum allows for ongoing adaptations based on feedback, ensuring its relevance and effectiveness.''(Yang, 2013) This adaptability paves the way for broader application across different educational settings, potentially benefiting a larger demographic of students.
  \item Empowerment Through Skill Development: ``Ultimately, the curriculum aims to empower students by equipping them with robust English writing skills, thereby enhancing their confidence and opening doors to global opportunities.'' (Yang, 2013)This empowerment is essential for students to actively participate in global discussions and pursue higher education and professional goals abroad.
\end{enumerate}

\section{Forward-Looking Statements}

The potential of this curriculum extends beyond immediate educational outcomes. As educators and stakeholders engage with this program, it is envisioned as a catalyst for broader educational reform, where the focus is on interactive, technology-driven learning across various subjects.

Looking ahead, the curriculum is set to evolve based on comprehensive feedback and the integration of advanced educational technologies. This evolution will aim to refine its effectiveness continually and expand its reach, ensuring it remains a cutting-edge tool that significantly contributes to global education initiatives. By doing so, the curriculum not only improves individual learning outcomes but also plays a pivotal role in fostering a more interconnected and understanding global community.

In conclusion, the Interactive English Writing Curriculum is more than just a teaching tool; it is a gateway to redefining educational paradigms and empowering the next generation of learners to succeed in a globally connected world. As we move forward, it remains a beacon of innovation and a testament to the transformative power of education.

\begin{references}
\item Graham, S., \& Perin, D. (2007). A meta-analysis of writing instruction for adolescent students. \textit{Journal of Educational Psychology, 99(3), 445–476.}
\item Yang, C. (2013). How Chinese Beginning Writers Learn English Writing: A Survey of Writing Strategies. \textit{Journal of Educational and Social Research}, \textit{3}(1). \url{https://doi.org/10.5901/jesr.2013.v3n1p9}
\item Wern, T. C., \& Rahmat, N. H. (2021). An investigative study on the types and causes of ESL writing anxiety: A case study of learners from a Chinese independent middle school. \textit{European Journal of English Language Teaching}, \textit{6}(3). \url{https://doi.org/10.46827/ejel.v6i3.3553}
\end{references}

\end{document}
